\section{系统的复频域分析}\label{sec:Complex_Freq_Analysis}

\subsection{系统的稳定性}\label{sub:stability_of_system}
本小节使用双边拉普拉斯变换分析系统的\textbf{稳定性}（\cref{def:system_stability}）。
\begin{proposition}[稳定性的时域判则]
    设线性时不变系统的冲激响应为$h(t)$，则系统稳定的充分必要条件是$\int_{-\infty}^{\infty}|h(t)|\,dt<\infty$.
\end{proposition}
\begin{proof}
    一方面，对于有界输入$|e(t)|<M$，有
    \begin{align*}
        |r(t)|&=\left|\int_{-\infty}^{\infty}h(\tau)e(t-\tau)\,d\tau\right|\\
        &\leq \int_{-\infty}^{\infty}|h(\tau)||e(t-\tau)|\,d\tau\\
        &<M\int_{-\infty}^{\infty}|h(\tau)|\,d\tau<\infty
    \end{align*}
    因此输出有界。

    另一方面，设$h(t)$不绝对可积，取有界激励信号$e(t)=\mathrm{sgn}(h(-t))$，则
    \begin{align*}
        |r(0)|&=\left|\int_{-\infty}^{\infty}h(\tau)e(-\tau)\,d\tau\right|\\
        &=\left|\int_{-\infty}^{\infty}h(\tau)\mathrm{sgn}(h(\tau))\,d\tau\right|\\
        &=\int_{-\infty}^{\infty}|h(\tau)|\,d\tau=\infty
    \end{align*}
    输出无界。
\end{proof}

显然，$\int_{-\infty}^{\infty}|h(t)|\,dt<\infty$与单边拉普拉斯变换的收敛域包含虚轴是等价的，由此可得：

\begin{proposition}[稳定性的复频域判则]\label{prop:statibility_s}
    设线性时不变系统的冲激响应为$h(t)\in\mathcal{E}$，系统函数为$H(s)=\mathcal{L} h(s)$，则稳定性的充分必要条件是$H(s)$的收敛域包含虚轴，即频率响应特性$H(\mathrm{i}\omega)$存在。
\end{proposition}

根据\refsub{converge_abscissa}，对于因果的线性时不变系统，这一判则在多数情况下还可以进一步简化为：$H(s)$的所有奇点均位于左半平面上，但严格来说它只是系统具有稳定性的必要条件，对于有理系统函数，它是充要条件。

\begin{definition}
    根据单位冲激响应的拉普拉斯变换的奇点位置，将因果线性时不变系统分为三类：
    \begin{itemize}
        \item \textbf{稳定系统}：所有奇点均位于左半平面上；
        \item \textbf{临界稳定系统}：至少有一个奇点位于虚轴上，且其余奇点均位于左半平面上；
        \item \textbf{不稳定系统}：至少有一个奇点位于右半平面上。
    \end{itemize}
\end{definition}


\subsection{最小相位系统}\label{sub:min_phase}

在上一小节中我们说明了：对于因果且稳定的线性时不变系统，系统函数$H(s)$的所有极点均位于复平面的左半平面。在设计模拟滤波器时，我们仅仅会给出因果、稳定的线性时不变系统\footnote{因果性是为了使滤波器可实现。至于佩利-维纳准则（\cref{def:paley_winner}），在幅度特性没有恒为0的区间时通常满足，因此很少单独验证它。}的幅度响应，如果我们还希望滤波器具有最小的时延特性，则需要令系统函数$H(s)$的所有零点也位于复平面的左半平面。

\begin{definition}[最小相位系统]
    设因果且稳定的连续线性时不变系统的系统函数为$H(s)$，记$H(s)$的零点集为$\{z_k\}$，如果$\forall k,\mathrm{Re}(z_k)<0$，则称该系统为\textbf{最小相位系统}(minimum-phase system)。
\end{definition}

为了说明为什么这样的系统在具有相同幅度响应的系统中具有最小的延时特性，我们首先考虑一个例子：
\begin{example}
    设因果且稳定的连续线性时不变系统A，B的系统函数分别为
    \[H_A(s)=\frac{s+1}{s+2},\quad H_B(s)=\frac{s-1}{s+2}\]
    这两个系统具有相同的幅频特性：
    \[|H_A(\mathrm{i}\omega)|=|H_B(\mathrm{i}\omega)|=\frac{\sqrt{\omega^2+1}}{\sqrt{\omega^2+4}}\]
    其中系统A为最小相位系统，系统B不是最小相位系统。计算两个系统的相频特性：
    \begin{align*}
        \varphi_A(\omega)&=\arg(\mathrm{i}\omega+1)-\arg(\mathrm{i}\omega+2)=\arctan(\omega)-\arctan\left(\frac{\omega}{2}\right)\\
        \varphi_B(\omega)&=\arg(\mathrm{i}\omega-1)-\arg(\mathrm{i}\omega+2)=-\arctan(\omega)-\arctan\left(\frac{\omega}{2}\right)
    \end{align*}

    系统的时延特由相时延$\tau_p$和群时延$\tau_g$（\cref{def:delay}）描述。计算两个系统的相时延和群时延：
    \begin{align*}
        \tau_{pA}(\omega)&=-\frac{\varphi_A(\omega)}{\omega}=\frac{-\arctan(\omega)+\arctan\left(\frac{\omega}{2}\right)}{\omega}\\
        \tau_{pB}(\omega)&=-\frac{\varphi_B(\omega)}{\omega}=\frac{\arctan(\omega)+\arctan\left(\frac{\omega}{2}\right)}{\omega}\\
        \tau_{gA}(\omega)&=-\frac{d\varphi_A(\omega)}{d\omega}=-\frac{1}{1+\omega^2}+\frac{2}{4+\omega^2}\\
        \tau_{gB}(\omega)&=-\frac{d\varphi_B(\omega)}{d\omega}=\frac{1}{1+\omega^2}+\frac{2}{4+\omega^2}
    \end{align*}
    可以发现，
    \begin{align*}
        &\tau_{pA}(\omega)-\tau_{pB}(\omega)=-\frac{2\arctan(\omega)}{\omega}<0\\
        &\tau_{gA}(\omega)-\tau_{gB}(\omega)=-\frac{2}{1+\omega^2}<0
    \end{align*}
    也就是说，最小相位系统A的相时延和群时延均小于非最小相位系统B的相时延和群时延，因此系统A具有更好的延时特性。
\end{example}

\begin{proposition}
    具有相同幅频特性的因果且稳定的连续线性时不变系统中，最小相位系统具有最小的相时延和群时延。
\end{proposition}
\begin{proof}
    我们对系统函数为有理函数的情况进行证明。设因果且稳定的连续线性时不变系统的系统函数为
    \[H(s)=\frac{K(s)\prod_{k=1}^{m}(s-z_k)}{\prod_{l=1}^{n}(s-p_l)}\]
    其中$K(s)$为实部小于0的零点对应因式的乘积，$z_k$为实部大于0的零点，$p_l$为实部小于0的极点。

    构造以下\textbf{全通系统}：
    $$A(s)=\frac{\prod_{k=1}^{m}(s-\tilde{z}_k)}{\prod_{k=1}^{m}(s-z_k)}$$
    其中$\tilde{z}_k=-\overline{z_k}$为$z_k$关于虚轴的对称点。设$z_k=\sigma_k+\mathrm{i}\omega_k(\sigma_k>0)$，则$\tilde{z}_k=-\sigma_k+\mathrm{i}\omega_k$，该系统的幅频特性恒为$1$：
    \begin{align*}
        |A(\mathrm{i}\omega)|=\frac{\prod_{k=1}^{m}\sqrt{(\omega-\omega_k)^2+\sigma_k^2}}{\prod_{k=1}^{m}\sqrt{(\omega-\omega_k)^2+\sigma_k^2}}=1
    \end{align*}

    将两系统级联，所得系统的系统函数为
    \[\tilde{H}(s)=H(s)A(s)=K(s)\frac{\prod_{k=1}^{m}(s-\tilde{z}_k)}{\prod_{l=1}^{n}(s-p_l)}\]
    也就是说，通过与一个全通系统级联，我们将系统原有的实部大于0的零点转化为实部小于0的零点，而保持$\tilde{H}(s)$的幅频特性与$H(s)$相同。此时，级联系统是最小相位系统。
    
    两个系统级联，相频特性相加：
    \begin{align*}
        \tilde{H}(s)=H(s)A(s),\quad \arg \tilde{H}(s)=\arg H(s)+\arg A(\omega)
    \end{align*}
    因此相时延和群时延也相加。为了验证级联系统具有最小相位性，我们来计算全通系统造成的相时延和群时延。全通系统的相频特性为：
    \begin{align*}
        \varphi_A(\omega)=\arg(A(\mathrm{i}\omega))=\sum_{k=1}^{m}\arg\left(\frac{\mathrm{i}\omega-\tilde{z}_k}{\mathrm{i}\omega-z_k}\right)=\sum_{k=1}^m \left[\arg\left(\sigma_k+\mathrm{i}(\omega-\omega_k)\right)-\arg\left(-\sigma_k+\mathrm{i}(\omega-\omega_k)\right)\right]
    \end{align*}
    选取解卷绕（见\refsub{distortion}）的辐角值，则
    \begin{align*}
        &\arg(\sigma_k+\mathrm{i}(\omega-\omega_k))-\arg(-\sigma_k+\mathrm{i}(\omega-\omega_k))\\
        =&\arctan\left(\frac{\omega-\omega_k}{\sigma_k}\right)-\left(\pi-\arctan\left(\frac{\omega-\omega_k}{\sigma_k}\right)\right)\\
        =&2\arctan\left(\frac{\omega-\omega_k}{\sigma_k}\right)-\pi
    \end{align*}
    因此\begin{align*}
        \varphi_{A,\mathrm{unwrapped}}(\omega)&=2\sum_{k=1}^m \arctan\left(\frac{\omega-\omega_k}{\sigma_k}\right)-m\pi
    \end{align*}

    计算全通系统的相时延和群时延：
    \begin{align*}
        \tau_{A,p}(\omega)&=-\frac{\varphi_{A,\mathrm{unwrapped}}(\omega)}{\omega}=\frac{1}{\omega}\sum_{k=1}^m \left[2\arctan\left(\frac{\omega-\omega_k}{\sigma_k}\right)-\pi\right]<0\\
        \tau_{A,g}(\omega)&=-\frac{d\varphi_{A,\mathrm{unwrapped}}(\omega)}{d\omega}=-2\sum_{k=1}^m \frac{\sigma_k}{\sigma_k^2+(\omega-\omega_k)^2}<0
    \end{align*}
    也就是说，全通系统的相时延和群时延均为负值。因此，级联系统的相时延和群时延均小于原系统。
\end{proof}

需要指出，这个证明并不严谨。其一，这个证明没有说明为什么最小相位系统要求系统函数在虚轴上也不存在零点；其二，这个证明仅适用于有理函数。但是，更一般的情况需要用到哈代空间中函数的内外分解，我没能找到既严谨又通俗的证明，只能说这个证明能够提供一些直观的理解。

\subsection{逆系统}\label{sub:inv_system}

\begin{definition}[逆系统]
    设因果的线性时不变系统的系统函数为$H(s)$，如果存在一个因果的连续线性时不变系统，其系统函数为$H^{-1}(s)$，则称该系统为原系统的\textbf{逆系统}。
\end{definition}

将原系统与其逆系统级联，系统函数为$H(s)H^{-1}(s)=1$，即恒等系统。

由于最小相位系统的系统函数$H(s)$的零、极点均位于复平面的左半平面，其倒数$H^{-1}(s)$的零、极点也均位于左半平面，因此最小相位系统的逆系统也是最小相位系统。

\begin{example}
    设因果且稳定的连续线性时不变系统的系统函数及单位脉冲响应为
    \[H(s)=\frac{s+1}{s+2},\quad h(t)=\delta(t)-\mathrm{e}^{-2t}u(t)\]
    则其逆系统的系统函数及单位脉冲响应为
    \[H^{-1}(s)=\frac{s+2}{s+1},\quad h_{\mathrm{inv}}(t)=\delta(t)+\mathrm{e}^{-t}u(t)\]
    该逆系统也是因果且稳定的连续线性时不变系统。将两个系统级联，其单位脉冲响应为
\begin{align*}
    h_{eq}(t)&=h(t)*h_{\mathrm{inv}}(t)\\
    &=\delta(t)-\mathrm{e}^{-2t}u(t)+\mathrm{e}^{-t}u(t)-(\mathrm{e}^{-2t}u(t))*(\mathrm{e}^{-t}u(t))\\
    &=\delta(t)-\mathrm{e}^{-2t}u(t)+\mathrm{e}^{-t}u(t)-u(t)\int_0^t \mathrm{e}^{-2\tau}\mathrm{e}^{-(t-\tau)}d\tau\\
    &=\delta(t)
\end{align*}
\end{example}

考虑原系统与逆系统级联时的相时延与群时延，由于两个系统的级联为恒等系统，它们的相时延之和与群时延之和必定为0，因此两个系统中必有其中之一具有负的时延。然而，\textbf{响应并不会先于激励出现，负的时延与因果性不冲突}。

为了明白这一点，我们首先考察上例中的情况。这个例子中，已经算出原系统的单位脉冲响应，它们都是因果信号，因此系统也是因果的。考虑它们的相频特性：
\begin{align*}
    \varphi(\omega)&=\arg(H(\mathrm{i}\omega))=\arctan(\omega)-\arctan\left(\frac{\omega}{2}\right)\\
    \varphi_{\mathrm{inv}}(\omega)&=\arg(H^{-1}(\mathrm{i}\omega))=\arctan\left(\frac{\omega}{2}\right)-\arctan(\omega)
\end{align*}
因此，相时延和群时延分别为
\begin{align*}
    \tau_p(\omega)&=-\frac{\varphi(\omega)}{\omega}=\frac{-\arctan(\omega)+\arctan\left(\frac{\omega}{2}\right)}{\omega}<0\\
    \tau_{p,\mathrm{inv}}(\omega)&=-\frac{\varphi_{\mathrm{inv}}(\omega)}{\omega}=\frac{-\arctan\left(\frac{\omega}{2}\right)+\arctan(\omega)}{\omega}>0\\
    \tau_g(\omega)&=-\frac{d\varphi(\omega)}{d\omega}=-\frac{1}{1+\omega^2}+\frac{2}{4+\omega^2}<0\\
    \tau_{g,\mathrm{inv}}(\omega)&=-\frac{d\varphi_{\mathrm{inv}}(\omega)}{d\omega}=-\frac{d(-\varphi(\omega))}{d\omega}=\frac{1}{1+\omega^2}-\frac{2}{4+\omega^2}>0
\end{align*}

那么，为什么负的延时不会使得系统违反因果性呢？实际上，这是因为单频正弦信号“提前出现”并不破坏因果性，因为当我们考虑系统的激励为正弦信号时，实际上假定这是一个\textbf{稳态}信号，即该正弦信号从$t=-\infty$时刻就已经存在了，我们观察到的相位“提前”，\textbf{实际上是相位延迟}，而级联所得的恒等变换系统总延迟为0，实际上是两个系统的相位延迟相互抵消了。因此，并不存在违反因果性的情况。

进一步，群时延之和为0，表现出调幅信号的包络峰提前出现的现象。这也没有违反因果性，因为响应起始时刻仍在激励之后，这种包络峰提前出现的现象，实际上是\textbf{相位失真引起的波形重构}。

\begin{example}[负的群时延]
    考虑调幅信号通过具有负的群时延的系统。设
    \[e(t)=\cos(\omega_m t)\cos(\omega_c t),\quad H(s)=\frac{s+1}{s+2},\quad h(t)=\delta-\mathrm{e}^{-2t}u(t)\]
    则系统的群时延为
    \[\tau_g(\omega)=-\frac{d\varphi(\omega)}{d\omega}=-\frac{1}{1+\omega^2}+\frac{2}{4+\omega^2}\]
    在低频段，群时延为负值。设基波频率$\omega_m=1$，载波频率$\omega_c=5$（这个载波的频率不算很高，这是为了使响应与激励区分开），则
    \begin{align*}
        e(t)=\cos(t)\cos(5t)=\frac{1}{2}[\cos(4t)+\cos(6t)],\quad \tau_g(1)=-0.1
    \end{align*}
    由此可以预测，系统的响应中，包络峰将提前出现0.1s。

    计算系统的响应：
    \begin{align*}
        y(t)&=(h*e)(t)\\
        &=\cos(t)\cos(5t)-\frac{1}{2}u(t)\int_{0}^{t}\mathrm{e}^{-2\tau}[\cos(4(t-\tau))+\cos(6(t-\tau))]d\tau\\
        &=\cos(t)\cos(5t)-\frac{1}{2}u(t)\left[\frac{\cos(4t)+2\sin(4t)-\mathrm{e}^{-2t}}{10}+\frac{\cos(6t)+3\sin(6t)-\mathrm{e}^{-2t}}{20}\right]
    \end{align*}
    \begin{figure}[H]
        \centering
        \includegraphics[width=0.8\textwidth]{neg_taug}
        \caption{调幅信号通过具有负群时延系统的响应}
    \end{figure}
    上图中，依稀能辨认出-0.1s的群时延，但由于输入的是稳态信号，并没有违反因果性。
\end{example}

那么，对于暂态信号，情况又如何呢？考虑相时延恒负的原系统，它会使得响应的各频率分量提前出现，响应的起始时刻可能会早于激励的起始时刻，这似乎违反了因果性。然而，响应的各频率分量的相时延不同，恰会使得响应在$t=0$时刻前抵消，最终呈现出响应与激励起始时刻相同，不违反因果性的结果。

\begin{example}[暂态信号经过原系统]
    设
    $$e(t)=\chi_{[0,1]}(t),\quad H(s)=\frac{s+1}{s+2}$$
    则系统的响应为
    \begin{align*}
        y(t)&=(h*\chi_{[0,1]})(t)\\
        &=\chi_{[0,1]}(t)-u(t)\int_0^t \chi_{[0,1]}(\tau)\mathrm{e}^{-2(t-\tau)}d\tau
    \end{align*}
    单位脉冲响应$h(t)$是因果的，因此响应也是因果的。

    在$t<1$时，相应为\begin{align*}
        y(t)&=\chi_{[0,1]}(t)-u(t)\int_0^t \mathrm{e}^{-2(t-\tau)}d\tau\\
        &=\chi_{[0,1]}(t)-\frac{1}{2}u(t)(1-\mathrm{e}^{-2t})
    \end{align*}
    在$t\geq 1$时，相应为\begin{align*}
        y(t)&=\chi_{[0,1]}(t)-u(t)\int_0^1 \mathrm{e}^{-2(t-\tau)}d\tau\\
        &=-\frac{1}{2}u(t)(\mathrm{e}^{-2t}-\mathrm{e}^{-2(t-1)})
    \end{align*}
    \begin{figure}[H]
        \centering
        \includegraphics[width=0.8\textwidth]{inv_system}
        \caption{暂态信号通过原系统的响应}
    \end{figure}
    响应与激励的起始时刻一致，满足因果性。
\end{example}

\subsection{零极点图与频率响应特性}\label{sub:zero_pole_characteristics}

对于最典型的因果线性时不变系统，其系统函数为有理函数，可以将系统函数表示为：
\begin{align*}
    H(s)=K\frac{(s-z_1)(s-z_2)\cdots (s-z_m)}{(s-p_1)(s-p_2)\cdots (s-p_n)}
\end{align*}
其中$z_1,z_2,\cdots ,z_m$为系统函数的零点，$p_1,p_2,\cdots ,p_n$为系统函数的极点，$m<n$，$K$为常数。将零点和极点在复平面上表示出来，得到\textbf{零极点图}(zero-pole plot)，其中，零点用“O”表示，极点用“X”表示，高阶的零点或极点，在其旁边用数字标明其重数。下面给出一个例子：

\begin{example}
    设系统函数为
    \[H(s)=\frac{s+2}{(s+1)^2 (s+3)}\]
    则其零极点图如下所示：
    \begin{figure}[H]
        \centering
        \includegraphics[width=0.4\textwidth]{zero_pole}
        \caption{零极点图(1)}
    \end{figure}
\end{example}

根据\cref{prop:statibility_s}，因果的线性时不变系统是稳定的，当且仅当其系统函数的收敛域包含虚轴，此时系统函数的所有极点均位于左半平面$\mathrm{Re}(s)<0$，\textbf{系统的频率响应特性即系统函数在虚轴上的限制}。

此时，如果零极点图较为简单，则可以由零极点图判断系统的滤波特性，即判断系统函数在虚轴上不同位置的幅值大小。

零点通常会使得附近的系统函数幅值减小，特别地，如果虚轴上存在零点，则可直接判断该点处的幅频特性值为0：

\begin{example}
    设系统的零极点图如下图(a)所示：
    \begin{figure}[H]
        \centering
        \includegraphics[width=0.8\textwidth]{zero_pole1}
        \caption{零极点图与幅频特性曲线(1)}
    \end{figure}
    可以看到系统函数在原点处由一个一阶零点，这表明系统的频率响应特性在$\omega=0$有零点，也就意味着系统对低频信号有抑制作用；系统的极点均位于实轴上，这意味着系统函数的分母是有两个实根的实系数多项式
    \[H(s)=\frac{Ks}{(s+a)(s+b)}\]
    考虑它在虚轴上的限制的幅值：
    \[|H(\mathrm{i}\omega)|=\frac{K\omega}{\sqrt{(\omega^2+a^2)(\omega^2+b^2)}}\]
    当$\omega\to\infty$时，$|H(\mathrm{i}\omega)|\to 0$，这表明系统对高频信号有抑制作用。因此，该系统具有带通滤波特性。

    可以定性地画出其幅频特性曲线，如上图(b)所示。
\end{example}

极点通常会使得附近的系统函数幅值增大：
\begin{example}
    设系统的零极点图如图下图(a)所示：
    \begin{figure}[H]
        \centering
        \includegraphics[width=0.7\textwidth]{zero_pole2}
        \caption{零极点图与幅频特性曲线(2)}
    \end{figure}

    系统函数的分母在虚轴左侧有一对共轭复根，这意味着系统函数形如
    \[H(s)=\frac{K}{(s-p)(s-\bar{p})}=\frac{K}{s^2-2\mathrm{Re}(p)s+|p|^2}\]
    考虑它在虚轴上的限制的幅值：
    \[|H(\mathrm{i}\omega)|=\frac{K}{\sqrt{\omega^2-2\mathrm{Re}(p)\omega+|p|^2}}\]
    由于$\mathrm{Re}(p)<0$，因此容易求得幅频特性在$\omega=\pm \mathrm{Im}(p)$处取得极大值，这表明系统对该频率附近的信号有增强作用。因此，该系统具有带通滤波特性。
    
    可以定性地画出其幅频特性曲线，如上图(b)所示。
\end{example}

以上推断系统函数的形式的过程，可以进一步简化。对于给定位置的零点$z$，其系统函数分子必然包含$s-z$这一因子，从而幅频特性是系统函数的其他部分与一个\textbf{起点为$z$，终点为$\mathrm{i}\omega$的向量长度}的积；对于给定位置的极点$p$，其系统函数分分母必然包含$s-p$这一因子，从而幅频特性是系统函数的其他部分与一个\textbf{起点为$p$，终点为$\mathrm{i}\omega$的向量长度}的商。

\begin{example}
    设系统的零极点图如下图(a)所示：
    \begin{figure}[H]
        \centering
        \includegraphics[width=0.8\textwidth]{zero_pole_vector}
        \caption{零极点图与幅频特性曲线(3)}
    \end{figure}

    考虑不同角频率处的各向量长度，当角频率较小时，各向量的长度都较小；随角频率增大，极点向量的长度增长较快，且极点有两个，因此系统的幅频特性大小随角频率增大而减小，系统具有低通滤波特性。
    
    可以定性地画出其幅频特性曲线，如上图(b)所示。虽然不太容易看出幅频特性曲线在低频段内有一个递增的区间，但这并不影响系统的低通滤波特性。
\end{example}